\clearpage
\chapter*{LAMPIRAN}
\addcontentsline{toc}{chapter}{LAMPIRAN}

\section*{Lampiran 1. Deskripsi Dataset \texttt{patient.csv}}
\label{app:patient}
\addcontentsline{toc}{section}{Lampiran 1. Deskripsi Dataset patient.csv}
\setlength{\tabcolsep}{3pt}
\begin{longtable}{|L{2.5cm}|C{1.8cm}|L{5.3cm}|L{3.5cm}|}
\hline
\textbf{Variabel} & \textbf{Tipe Data} & \textbf{Deskripsi} & \textbf{Keterangan}\\ \hline
\endfirsthead
\hline
\textbf{Variabel} & \textbf{Tipe Data} & \textbf{Deskripsi} & \textbf{Keterangan}\\ \hline
\endhead
\texttt{PatientGuid} & String & Identitas unik untuk setiap pasien yang digunakan sebagai kunci utama dalam proses penggabungan data. & ID alfanumerik.\\ \hline
\texttt{Age} & Integer & Usia pasien pada saat pencatatan data. & 0--120 tahun.\\ \hline
\texttt{Gender} & Kategorikal & Jenis kelamin pasien. & Male, Female, Unknown.\\ \hline
\texttt{State} & Kategorikal & Lokasi geografis pasien berdasarkan wilayah negara bagian. & Kode state (CA, TX, NY, dan lain-lain).\\ \hline
\texttt{DMIndicator} & Biner (Target) & Status diagnosis Diabetes Mellitus Tipe 2. & 1 = Diabetes, 0 = Non-diabetes.\\ \hline
\end{longtable}

\section*{Lampiran 2. Deskripsi Dataset \texttt{diagnosis.csv}}
\label{app:diagnosis}
\addcontentsline{toc}{section}{Lampiran 2. Deskripsi Dataset diagnosis.csv}
\small
\begin{longtable}{|L{2.7cm}|C{1.8cm}|L{5.1cm}|L{3.5cm}|}
\hline
\textbf{Variabel} & \textbf{Tipe Data} & \textbf{Deskripsi} & \textbf{Keterangan}\\ \hline
\endfirsthead
\hline
\textbf{Variabel} & \textbf{Tipe Data} & \textbf{Deskripsi} & \textbf{Keterangan}\\ \hline
\endhead
\texttt{PatientGuid} & String & Identitas unik pasien. & Digunakan sebagai kunci penggabungan data.\\ \hline
\texttt{Icd9\_001-139} & Biner (0/1) & Penyakit infeksi dan parasit (rentang ICD-9: 001--139). & Menunjukkan keberadaan diagnosis.\\ \hline
\texttt{Icd9\_140-239} & Biner (0/1) & Neoplasma (rentang ICD-9: 140--239). & Riwayat kanker atau tumor.\\ \hline
\texttt{Icd9\_240-279} & Biner (0/1) & Penyakit endokrin, nutrisi, dan metabolik. & Sangat relevan terhadap diabetes.\\ \hline
\texttt{Icd9\_280-289} & Biner (0/1) & Penyakit darah dan organ pembentuk darah. & Misalnya anemia.\\ \hline
\texttt{Icd9\_290-319} & Biner (0/1) & Gangguan mental, perilaku, dan perkembangan saraf. & Riwayat gangguan psikiatri.\\ \hline
\texttt{Icd9\_320-359} & Biner (0/1) & Penyakit sistem saraf. & Kondisi neurologis.\\ \hline
\texttt{Icd9\_360-389} & Biner (0/1) & Penyakit mata dan telinga. & Gangguan sistem sensorik.\\ \hline
\texttt{Icd9\_390-459} & Biner (0/1) & Penyakit sistem peredaran darah. & Berkaitan dengan hipertensi dan penyakit jantung.\\ \hline
\texttt{Icd9\_460-519} & Biner (0/1) & Penyakit sistem pernapasan. & Gangguan paru-paru.\\ \hline
\texttt{Icd9\_520-579} & Biner (0/1) & Penyakit sistem pencernaan. & Gangguan gastrointestinal.\\ \hline
\texttt{Icd9\_580-629} & Biner (0/1) & Penyakit sistem genitourinari. & Termasuk komplikasi ginjal.\\ \hline
\texttt{Icd9\_630-679} & Biner (0/1) & Komplikasi kehamilan, persalinan, dan masa nifas. & Kondisi obstetri.\\ \hline
\texttt{Icd9\_680-709} & Biner (0/1) & Penyakit kulit dan jaringan subkutan. & Gangguan kulit.\\ \hline
\texttt{Icd9\_710-739} & Biner (0/1) & Penyakit sistem muskuloskeletal. & Gangguan tulang dan sendi.\\ \hline
\texttt{Icd9\_740-759} & Biner (0/1) & Kelainan bawaan. & Cacat sejak lahir.\\ \hline
\texttt{Icd9\_760-779} & Biner (0/1) & Kondisi yang berasal dari periode perinatal. & Kondisi terkait kelahiran.\\ \hline
\texttt{Icd9\_780-799} & Biner (0/1) & Gejala, tanda, dan temuan abnormal. & Gejala non-spesifik.\\ \hline
\texttt{Icd9\_800-899} & Biner (0/1) & Cedera dan keracunan akibat faktor eksternal. & Riwayat trauma.\\ \hline
\texttt{Icd9\_E-V} & Biner (0/1) & Penyebab eksternal cedera dan klasifikasi tambahan. & Kode E dan V.\\ \hline
\texttt{Icd9\_Unknown} & Biner (0/1) & Kode diagnosis tidak diketahui atau tidak terklasifikasi. & Indikator kualitas data.\\ \hline
\texttt{DiagnosisCount} & Integer & Jumlah total diagnosis yang dimiliki pasien. & Menggambarkan beban penyakit.\\ \hline
\texttt{VisitCount} & Integer & Jumlah total kunjungan medis pasien. & Tingkat interaksi dengan layanan kesehatan.\\ \hline
\texttt{DiagnosisFreq} & Float & Rata-rata jumlah diagnosis per kunjungan. & Indikator kompleksitas klinis.\\ \hline
\texttt{AcuteCount} & Integer & Jumlah diagnosis akut (kondisi sementara). & Frekuensi penyakit akut.\\ \hline
\texttt{AcuteFreq} & Float & Rata-rata diagnosis akut per kunjungan. & Tingkat kejadian kondisi akut.\\ \hline
\end{longtable}
\normalsize

\section*{Lampiran 3. Deskripsi Dataset \texttt{medication.csv}}
\label{app:medication}
\addcontentsline{toc}{section}{Lampiran 3. Deskripsi Dataset medication.csv}
\small
\begin{longtable}{|L{2.8cm}|C{1.8cm}|L{3.3cm}|L{5.2cm}|}
\hline
\textbf{Variabel} & \textbf{Tipe Data} & \textbf{Deskripsi} & \textbf{Keterangan}\\ \hline
\endfirsthead
\hline
\textbf{Variabel} & \textbf{Tipe Data} & \textbf{Deskripsi} & \textbf{Keterangan}\\ \hline
\endhead
\texttt{PatientGuid} & String & Identitas unik pasien. & Digunakan sebagai kunci penggabungan data.\\ \hline
\texttt{Medication\_1} hingga \texttt{Medication\_n} & Biner (0/1) & Indikator apakah suatu obat pernah diresepkan kepada pasien. & Contoh: abacavir, abilify, acarbose, aceon, aciphex, acyclovir, adalimumab, adderall, advair, advil, albuterol, aldactone, aleve, allegra, allopurinol, amaryl, ambien, amiodarone, amitriptyline, amlodipine, amoxicillin, ampicillin, anastrozole, androgel, aspirin, atenolol, atorvastatin, azithromycin, furosemide, hydrochlorothiazide, ibuprofen, insulin, lantus, lisinopril, losartan, metformin, metoprolol, naproxen, omeprazole, pravastatin, simvastatin, spironolactone, dan lain-lain.\\ \hline
\texttt{DMIndicator} & Biner & Status diagnosis Diabetes Mellitus Tipe 2. & 1 = Diabetes, 0 = Non-diabetes.\\ \hline
\end{longtable}
\normalsize

\section*{Lampiran 4. Deskripsi Dataset \texttt{transcript.csv}}
\label{app:transcript}
\addcontentsline{toc}{section}{Lampiran 4. Deskripsi Dataset transcript.csv}
\scriptsize
\begin{longtable}{|L{2.8cm}|L{4.2cm}|L{3.5cm}|L{2.6cm}|}
\hline
\textbf{Kategori} & \textbf{Variabel} & \textbf{Deskripsi} & \textbf{Relevansi terhadap Diabetes}\\ \hline
\endfirsthead
\hline
\textbf{Kategori} & \textbf{Variabel} & \textbf{Deskripsi} & \textbf{Relevansi terhadap Diabetes}\\ \hline
\endhead
Anthropometric & \path{Height_Max, Height_Min, Height_Mean, Height_Std} & Tinggi badan pasien dalam berbagai nilai agregasi. & Digunakan sebagai dasar perhitungan BMI.\\ \hline
Anthropometric & \path{Weight_Max, Weight_Min, Weight_Mean, Weight_Std} & Berat badan pasien dalam berbagai nilai agregasi. & Faktor penting dalam risiko T2DM.\\ \hline
Anthropometric & \path{BMI_Max, BMI_Min, BMI_Mean, BMI_Std} & Indeks Massa Tubuh yang dihitung dari berat dan tinggi badan. & Prediktor utama risiko diabetes.\\ \hline
Anthropometric & \path{Weight_Change} & Perubahan berat badan antarkunjungan. & Indikator tren kondisi pasien.\\ \hline
Blood Pressure & \path{SystolicBP_Max, SystolicBP_Min, SystolicBP_Mean, SystolicBP_Std} & Tekanan darah sistolik dalam berbagai agregasi. & Berkaitan dengan risiko hipertensi.\\ \hline
Blood Pressure & \path{DiastolicBP_Max, DiastolicBP_Min, DiastolicBP_Mean, DiastolicBP_Std} & Tekanan darah diastolik dalam berbagai agregasi. & Indikator tambahan kondisi kardiovaskular.\\ \hline
Blood Pressure & \path{SystolicBP_Change, DiastolicBP_Change} & Perubahan tekanan darah antarwaktu. & Indikator perkembangan kondisi kesehatan.\\ \hline
Respiratory & \path{RespiratoryRate_Max, RespiratoryRate_Min, RespiratoryRate_Mean, RespiratoryRate_Std} & Laju pernapasan pasien. & Indikator kondisi kesehatan umum.\\ \hline
Respiratory & \path{RespiratoryRate_Change} & Perubahan laju pernapasan. & Indikator penurunan kondisi kesehatan.\\ \hline
Thermoregulation & \path{Temperature_Max, Temperature_Min, Temperature_Mean, Temperature_Std} & Suhu tubuh pasien. & Indikator infeksi atau inflamasi.\\ \hline
Thermoregulation & \path{Temperature_Change} & Perubahan suhu tubuh antarwaktu. & Indikator perubahan kondisi kesehatan.\\ \hline
Derived Metrics & \path{BMI_Change} & Perubahan BMI dari waktu ke waktu. & Penting untuk analisis tren risiko diabetes.\\ \hline
Derived Metrics & \path{Height_Change} & Perubahan tinggi badan. & Digunakan untuk pengecekan kualitas data.\\ \hline
\end{longtable}
\normalsize

\setlength{\tabcolsep}{6pt}
